% ===========================================================================
% Paper 51: Prime Gaps, Twin Primes and the Andrica-Cramér Conjecture
% Author: Michael Fuhrmann
% Project: Specialist Mathematics Project
% Build: 137
% Date: 2026-03-12
% ===========================================================================
\documentclass[12pt,a4paper]{amsart}
\usepackage{amsmath,amssymb,amsthm}
\usepackage{mathtools}
\usepackage[utf8]{inputenc}
\usepackage[T1]{fontenc}
\usepackage{hyperref}
\usepackage{enumitem,booktabs,array}
\usepackage{geometry}
\geometry{margin=2.5cm}

\newtheorem{theorem}{Theorem}[section]
\newtheorem{lemma}[theorem]{Lemma}
\newtheorem{corollary}[theorem]{Corollary}
\newtheorem{proposition}[theorem]{Proposition}
\newtheorem{conjecture}[theorem]{Conjecture}
\theoremstyle{definition}
\newtheorem{definition}[theorem]{Definition}
\newtheorem{example}[theorem]{Example}
\newtheorem{remark}[theorem]{Remark}

\author{Michael Fuhrmann}
\address{Specialist Mathematics Project, Build 137}
\email{specialist-maths@localhost}
\date{\today}

\title{Prime Gaps, Twin Primes, and the Andrica--Cram\'{e}r Conjecture:\\
A Survey of Open Problems and Modern Breakthroughs}

\begin{document}

\begin{abstract}
The distribution of prime gaps---the differences $g_n = p_{n+1} - p_n$ between
consecutive primes---lies at the heart of analytic number theory. This paper
surveys the principal conjectures and proved results concerning prime gaps, twin
primes, and related problems. We treat the celebrated Twin Prime Conjecture
(open since antiquity), the Andrica Conjecture ($\sqrt{p_{n+1}} - \sqrt{p_n} < 1$,
open), the Cramér Conjecture ($\limsup g_n/(\log p_n)^2 = 1$, open), Legendre's
Conjecture (open), and Polignac's Conjecture (open for all specific values).
The landmark breakthroughs of Zhang (2013, $\liminf g_n \le 70{,}000{,}000$,
\textbf{proved}) and Maynard--Polymath8b (2014, $\liminf g_n \le 246$,
\textbf{proved}) are presented in full, together with the GPY sieve technique
and Brun's theorem on the Brun constant. Probabilistic models (Cramér, Granville)
and extensive computational evidence are discussed. We carefully distinguish all
proved theorems from open conjectures. Numerical experiments were performed
using the Python modules \texttt{twin\_prime\_analysis.py} and
\texttt{andrica\_cramer.py} from the Specialist Mathematics Project (Build 137).
\end{abstract}

\maketitle
\tableofcontents

% ---------------------------------------------------------------------------
\section{Introduction}
% ---------------------------------------------------------------------------

Prime numbers are the multiplicative atoms of the integers, yet their additive
distribution conceals deep and largely unsolved mysteries. Among the most
fundamental questions in analytic number theory is: \emph{How large can the gap
between consecutive primes become, and how often do primes occur very close
together?}

The \emph{prime gap} $g_n = p_{n+1} - p_n$ between the $n$-th prime $p_n$ and
its successor $p_{n+1}$ has been studied since antiquity. Euclid's proof of the
infinitude of primes, combined with the Chinese Remainder Theorem, shows that
$g_n$ can be arbitrarily large. By contrast, whether $g_n = 2$ infinitely
often---the celebrated \emph{Twin Prime Conjecture}---remains one of the most
famous open problems in all of mathematics.

The connection to the Riemann Hypothesis (RH) is direct: under the assumption of
RH, one obtains the conditional bound
\[
  g_n \ll \sqrt{p_n} \log p_n,
\]
which is far better than any unconditional result currently available. The
unresolved nature of both RH and the Twin Prime Conjecture illustrates that prime
gaps sit at the intersection of the deepest open questions in number theory.

This paper surveys the principal conjectures and proven results concerning prime
gaps. We distinguish carefully between \textbf{proven theorems} (marked as
Theorem) and \textbf{open conjectures} (marked as Conjecture). The landmark
results of Zhang (2013) and Maynard (2014), which established that
$\liminf_{n\to\infty} g_n < \infty$, are treated in full. The Andrica
Conjecture ($\sqrt{p_{n+1}} - \sqrt{p_n} < 1$), the Cram\'{e}r Conjecture
($\limsup g_n / (\log p_n)^2 = 1$), Legendre's Conjecture, and Polignac's
Conjecture are presented with their current status, computational evidence, and
mutual implications.

All numerical computations were performed using the Python modules
\texttt{twin\_prime\_analysis.py} and \texttt{andrica\_cramer.py} from the
Specialist Mathematics Project (Build 137).

% ---------------------------------------------------------------------------
\section{Prime Gaps}
% ---------------------------------------------------------------------------

\begin{definition}
Let $p_1 = 2 < p_2 = 3 < p_3 = 5 < \cdots$ be the sequence of prime numbers in
increasing order. The \emph{$n$-th prime gap} is
\[
  g_n := p_{n+1} - p_n, \quad n \geq 1.
\]
A \emph{maximal gap} (also called \emph{record gap} or \emph{prime desert}) is a
gap $g_n$ that exceeds all preceding gaps: $g_n > g_k$ for all $k < n$.
\end{definition}

\begin{example}
The first few prime gaps are $g_1 = 1$ (from $2$ to $3$), then
$g_2 = 2, g_3 = 2, g_4 = 4, g_5 = 2, g_6 = 4, g_7 = 2, g_8 = 4, g_9 = 6,
\ldots$ All gaps $g_n \geq 2$ for $n \geq 2$ since all primes $p \geq 3$ are odd.
\end{example}

\begin{theorem}[Infinitude of arbitrarily large gaps]
For every integer $k \geq 1$, there exists a prime gap $g_n \geq k$.
\end{theorem}

\begin{proof}
Consider $m = (k+1)! + 2$. The integers $m, m+1, \ldots, m+k-1$ are all
composite: $m+j = (k+1)! + (j+2)$ is divisible by $j+2$ for $0 \leq j \leq
k-1$. Hence no prime lies in this interval of length $k-1$, giving a gap of
size at least $k$.
\end{proof}

\subsection{Average Gap Size and the Prime Number Theorem}

The \emph{Prime Number Theorem} (PNT), proved independently by Hadamard and de la
Vall\'{e}e Poussin in 1896, states
\[
  \pi(x) \sim \frac{x}{\log x}, \quad x \to \infty,
\]
where $\pi(x)$ counts the number of primes $\leq x$. An equivalent formulation
is $p_n \sim n \log n$. It follows that the \emph{average} gap near $p_n$ is
approximately $\log p_n$:
\[
  \frac{1}{N} \sum_{n=1}^{N} g_n = \frac{p_{N+1} - 2}{N} \sim \log p_N.
\]

\subsection{Cram\'{e}r's Probabilistic Model}

Harald Cram\'{e}r (1936) introduced a powerful heuristic: model the primes as a
random subset of the integers where each integer $n$ is independently ``prime''
with probability $1/\log n$. Under this model:
\begin{itemize}
  \item The probability that both $n$ and $n + 2$ are prime is
        $\approx 1/(\log n)^2$.
  \item The expected number of primes in $[N, N + (\log N)^2]$ is $O(1)$.
  \item Maximal gaps up to $x$ satisfy $g_{\max} \approx (\log x)^2$.
\end{itemize}

\begin{remark}
The Cram\'{e}r model does not account for divisibility by small primes (the
``sieve effect''). Granville (1995) showed this leads to a corrected prediction:
the limsup of $g_n/(\log p_n)^2$ should be $2e^{-\gamma} \approx 1.1229$ rather
than $1$, where $\gamma \approx 0.5772$ is the Euler--Mascheroni constant.
\end{remark}

\subsection{Known Unconditional Bounds on Gaps}

Baker, Harman, and Pintz (2001) proved unconditionally:
\[
  g_n \ll p_n^{0.525}.
\]
Under RH, the bound improves to $g_n \ll p_n^{1/2} \log p_n$.

The Erd\H{o}s--Rankin construction gives a lower bound for maximal gaps:
there exist infinitely many $n$ with
\[
  g_n \gg \frac{\log p_n \cdot \log\log p_n \cdot \log\log\log\log p_n}{(\log\log\log p_n)^2}.
\]

% ---------------------------------------------------------------------------
\section{The Twin Prime Conjecture}
% ---------------------------------------------------------------------------

\begin{definition}
A \emph{twin prime pair} is a pair $(p, p+2)$ of prime numbers differing by $2$.
The \emph{twin prime counting function} is
\[
  \pi_2(x) := \#\{p \leq x : p \text{ and } p+2 \text{ both prime}\}.
\]
\end{definition}

The first twin prime pairs are: $(3,5), (5,7), (11,13), (17,19), (29,31),
(41,43), (59,61), (71,73), \ldots$

\begin{conjecture}[Twin Prime Conjecture, folklore; open since antiquity]
There are infinitely many prime pairs $(p, p+2)$. Equivalently,
$\pi_2(x) \to \infty$ as $x \to \infty$.
\end{conjecture}

\textbf{Status: Open. No proof is known.} The conjecture has been verified
computationally for all primes up to $10^{15}$.

\subsection{Brun's Theorem and the Brun Constant}

\begin{theorem}[Brun, 1919]
The series
\[
  B_2 := \sum_{\substack{p,\, p+2 \text{ twin primes}}} \left(\frac{1}{p} + \frac{1}{p+2}\right)
\]
converges. The \emph{Brun constant} is estimated as $B_2 \approx 1.9021605831$.
\end{theorem}

This result, proved by Viggo Brun using sieve methods, is remarkable: it holds
regardless of whether there are finitely or infinitely many twin primes. If the
twin primes were infinite but sparse enough, convergence would still hold.

\begin{remark}
Compare with the \emph{sum of reciprocals of all primes}:
$\sum_p 1/p = \infty$ (Euler). The convergence of $B_2$ contrasts sharply
and shows that twin primes (if infinite) are much sparser than primes themselves.
\end{remark}

\subsection{Hardy--Littlewood Conjecture B}

\begin{conjecture}[Hardy--Littlewood Conjecture B, 1923; open]
\label{conj:hl}
As $x \to \infty$,
\[
  \pi_2(x) \sim 2 C_2 \frac{x}{(\log x)^2},
\]
where $C_2$ is the \emph{twin prime constant}:
\[
  C_2 = \prod_{\substack{p \geq 3 \\ p \text{ prime}}} \frac{p(p-2)}{(p-1)^2}
       \approx 0.6601618158\ldots
\]
\end{conjecture}

The constant $C_2$ arises from the density correction factor for the sieve of
Eratosthenes applied simultaneously to $n$ and $n+2$. Numerically,
$2C_2 \approx 1.3203$.

\begin{remark}
This conjecture implies the Twin Prime Conjecture but is itself much stronger.
It has been verified to be consistent with numerical data up to $10^{15}$.
A proof seems far beyond current techniques.
\end{remark}

\subsection{The Parity Problem}

A fundamental obstruction to proving the Twin Prime Conjecture via classical
sieve methods is the \emph{Selberg parity problem}. Let $\Omega(n)$ denote the
number of prime factors of $n$ counted with multiplicity. Classical sieves
cannot distinguish the parity of $\Omega(n)$, making them unable to directly
count pairs where both $n$ and $n+2$ are prime (i.e., $\Omega(n) =
\Omega(n+2) = 1$). This is because any sieve function $\lambda^2$ treats both
the case $\Omega(n) = 1$ and $\Omega(n) = 2$ identically.

% ---------------------------------------------------------------------------
\section{Zhang's Theorem and Maynard's Improvement}
% ---------------------------------------------------------------------------

The period 2013--2015 witnessed one of the most dramatic breakthroughs in
analytic number theory in decades.

\begin{theorem}[Zhang, 2013]
\label{thm:zhang}
There exist infinitely many pairs of consecutive primes $(p_n, p_{n+1})$ with
\[
  p_{n+1} - p_n < 70{,}000{,}000.
\]
Equivalently, $\liminf_{n \to \infty} g_n \leq 70{,}000{,}000$.
\end{theorem}

Zhang's proof marked the first time anyone had shown that $\liminf g_n < \infty$.
The key ingredients were:
\begin{enumerate}[label=(\roman*)]
  \item The \emph{Goldston--Pintz--Y{\i}ld{\i}r{\i}m (GPY) sieve} (2005), which
        had come tantalizingly close but could not finish the argument.
  \item A strengthened estimate for primes in arithmetic progressions, using
        exponential sum techniques beyond the reach of the Bombieri--Vinogradov
        theorem.
  \item An admissible $k$-tuple of size $k = 3.5 \times 10^6$ whose elements
        collectively guaranteed a bounded gap.
\end{enumerate}

\begin{theorem}[Maynard, 2014; Polymath8b, 2014]
\label{thm:maynard}
Under the Goldston--Pintz--Y{\i}ld{\i}r{\i}m framework with improved weights,
\[
  \liminf_{n \to \infty} g_n \leq 246.
\]
More precisely, the Polymath8b project (a collaborative optimization of Maynard's
method) established this bound.
\end{theorem}

Maynard introduced a multi-dimensional sieve weight of the form
\[
  w(n) = \left(\sum_{d_1 | (n+h_1),\, d_2 | (n+h_2),\, \ldots} \lambda_{d_1, \ldots, d_k}\right)^2
\]
and optimized over the $\lambda$ coefficients to show that for any admissible
$k$-tuple $(h_1,\ldots,h_k)$ with $k$ large enough, at least two of the $n+h_i$
are prime infinitely often.

\begin{remark}
The bound $246$ means: there exists an even integer $2 \leq d \leq 246$ such
that $p_{n+1} - p_n = d$ for infinitely many $n$. The precise value $d = 2$
(twin primes) remains unproved.
\end{remark}

\subsection{The GPY Sieve Technique}

The GPY approach (Goldston, Pintz, Y{\i}ld{\i}r{\i}m, 2005) established:
\begin{theorem}[GPY, 2005]
For every $\varepsilon > 0$, there exist infinitely many $n$ with
\[
  g_n < \varepsilon \log p_n.
\]
\end{theorem}

This says primes sometimes cluster much closer than average, but does not give a
fixed bound. The key idea: introduce a \emph{mollifier} weight $w(n) = (\sum_{d |
P(n)} \lambda_d)^2$ where $P(n) = n(n+2)$ and optimize $\lambda_d$ to maximize
\[
  S = \frac{\sum_{n \sim N} w(n)[\Lambda(n+h_1)\Lambda(n+h_2)]}{\sum_{n \sim N} w(n)}.
\]

% ---------------------------------------------------------------------------
\section{Andrica's Conjecture}
% ---------------------------------------------------------------------------

\begin{conjecture}[Andrica, 1985; open]
\label{conj:andrica}
For all $n \geq 1$,
\[
  A_n := \sqrt{p_{n+1}} - \sqrt{p_n} < 1.
\]
\end{conjecture}

\textbf{Status: Open. Verified numerically for all primes up to $4 \times
10^{18}$.}

\subsection{Equivalent Formulations}

The Andrica Conjecture is equivalent to each of the following:
\begin{enumerate}[label=(\roman*)]
  \item $\sqrt{p_{n+1}} - \sqrt{p_n} < 1$ (original form),
  \item $p_{n+1} < (\sqrt{p_n} + 1)^2 = p_n + 2\sqrt{p_n} + 1$,
  \item $g_n = p_{n+1} - p_n < 2\sqrt{p_n} + 1$.
\end{enumerate}

Formulation (iii) makes the content transparent: the $n$-th prime gap must be
sublinear in $\sqrt{p_n}$.

\subsection{The Maximal Andrica Value}

The largest known value of $A_n$ occurs at $n = 4$:
\[
  A_4 = \sqrt{11} - \sqrt{7} \approx 3.3166 - 2.6457 = 0.6708.
\]
No value $A_n \geq 1$ has ever been found. As $n \to \infty$, $A_n$ appears to
decrease toward $0$, consistent with the expectation that prime gaps grow much
slower than $2\sqrt{p_n}$.

\begin{table}[h]
\centering
\caption{First Andrica values $A_n = \sqrt{p_{n+1}} - \sqrt{p_n}$}
\begin{tabular}{@{}rrrrc@{}}
\toprule
$n$ & $p_n$ & $p_{n+1}$ & $g_n$ & $A_n$ \\
\midrule
1 & 2 & 3 & 1 & 0.3178 \\
2 & 3 & 5 & 2 & 0.5058 \\
3 & 5 & 7 & 2 & 0.3780 \\
4 & 7 & 11 & 4 & \textbf{0.6708} \\
5 & 11 & 13 & 2 & 0.2970 \\
6 & 13 & 17 & 4 & 0.5353 \\
7 & 17 & 19 & 2 & 0.2375 \\
8 & 19 & 23 & 4 & 0.4508 \\
9 & 23 & 29 & 6 & 0.6065 \\
10 & 29 & 31 & 2 & 0.1826 \\
\bottomrule
\end{tabular}
\end{table}

\subsection{Relation to Other Conjectures}

\begin{proposition}
The Cram\'{e}r Conjecture implies the Andrica Conjecture.
\end{proposition}

\begin{proof}
Under the Cram\'{e}r Conjecture, $g_n \leq (1+\varepsilon)(\log p_n)^2$ for all
sufficiently large $n$. For $p$ large enough, $(\log p)^2 \ll \sqrt{p}$, so
$g_n \ll \sqrt{p_n}$. Hence $g_n < 2\sqrt{p_n} + 1$ for large $n$, and the
finitely many small $n$ can be checked directly.
\end{proof}

\begin{proposition}
The Andrica Conjecture implies the Legendre Conjecture for sufficiently large $n$.
\end{proposition}

The connection follows from equivalent formulation (ii) above.

% ---------------------------------------------------------------------------
\section{Cram\'{e}r's Conjecture}
% ---------------------------------------------------------------------------

\begin{conjecture}[Cram\'{e}r, 1936; open]
\label{conj:cramer}
\[
  \limsup_{n \to \infty} \frac{g_n}{(\log p_n)^2} = 1.
\]
\end{conjecture}

\textbf{Status: Open. The limsup is known to be at least $2e^{-\gamma} \approx
1.1229$ under heuristic arguments (Granville 1995), and is finite by
Baker--Harman--Pintz.}

\subsection{Granville's Correction}

Andrew Granville (1995) argued that the correct prediction from a more careful
analysis of the Cram\'{e}r model---incorporating Hardy--Littlewood
corrections---should be:
\[
  \limsup_{n \to \infty} \frac{g_n}{(\log p_n)^2} = 2e^{-\gamma} \approx 1.1229,
\]
where $\gamma = 0.5772156649\ldots$ is the Euler--Mascheroni constant.

The Granville correction $C_{\text{Gr}} = 2e^{-\gamma} \approx 1.1229$ arises
because the Cram\'{e}r model slightly overestimates the density of primes near
smooth numbers (numbers with only small prime factors), and this bias accumulates.

\subsection{Known Bounds}

\begin{theorem}[Baker--Harman--Pintz, 2001]
Unconditionally, $g_n \ll p_n^{0.525}$.
\end{theorem}

\begin{theorem}[Conditional on RH]
$g_n \ll p_n^{1/2} \log p_n$.
\end{theorem}

\begin{theorem}[Westzynthius, 1931]
$\limsup_{n\to\infty} g_n/\log p_n = \infty$.
\end{theorem}

\subsection{Record Cram\'{e}r Values}

The largest known normalized gap $g_n/(\log p_n)^2$ was found by Oliveira e
Silva (2014) at $p_n = 1\,693\,182\,318\,746\,371 \approx 1.7 \times 10^{15}$,
with normalized value $\approx 1.17$---well below the Granville bound of
$\approx 1.1229 \times (\log p)^2$ but consistent with it.

% ---------------------------------------------------------------------------
\section{Legendre's Conjecture}
% ---------------------------------------------------------------------------

\begin{conjecture}[Legendre, $\sim$1798; open]
\label{conj:legendre}
For every positive integer $n \geq 1$, there exists at least one prime $p$ with
\[
  n^2 < p < (n+1)^2.
\]
\end{conjecture}

\textbf{Status: Open. Verified computationally for all $n \leq 10^{15}$
(Meissel--Lehmer methods).}

\begin{remark}
The Legendre Conjecture is equivalent to saying that consecutive perfect squares
always contain a prime between them. This would imply that the gap between
consecutive primes near $n^2$ satisfies $g \ll n \sim \sqrt{p}$, which is
weaker than the Cram\'{e}r Conjecture but still unproved.
\end{remark}

\subsection{Connection to Andrica's Conjecture}

\begin{proposition}
The Andrica Conjecture is stronger than Legendre's Conjecture in the following
sense: if $A_n < 1$ for all $n$, then in every interval $(m^2, (m+1)^2)$ there
is at most one ``gap'' from the Andrica sequence.
\end{proposition}

The precise relationship is:

\begin{itemize}
  \item Andrica for consecutive primes: $p_{n+1} < (\sqrt{p_n}+1)^2$.
  \item Legendre: In every $(n^2, (n+1)^2)$ there is a prime.
\end{itemize}

These are not identical because Legendre considers all squares $n^2$, while
Andrica considers only pairs of consecutive primes.

\subsection{Partial Results}

Ingham (1937) proved that there is a prime in $(n^3, (n+1)^3)$ for all
sufficiently large $n$. Under RH, one can show there is a prime in
$(x, x + C\sqrt{x}\log x)$ for an explicit constant $C$.

% ---------------------------------------------------------------------------
\section{Polignac's Conjecture}
% ---------------------------------------------------------------------------

\begin{conjecture}[Polignac, 1849; open]
\label{conj:polignac}
For every positive even integer $2k$ ($k \geq 1$), there are infinitely many
prime pairs $(p, p + 2k)$.
\end{conjecture}

\textbf{Status: Open for all specific values of $k$, including $k=1$ (Twin Prime
Conjecture).}

The Twin Prime Conjecture is the special case $k=1$.

\begin{remark}
Thanks to Maynard's Theorem (\ref{thm:maynard}), we know that \emph{some}
even integer $2k \leq 246$ satisfies Polignac's conjecture. But for no specific
value of $k$ has this been proved.
\end{remark}

\subsection{de Polignac's Original Claim}

Alphonse de Polignac (1849) actually claimed a stronger result---that every even
number $2k$ is the difference of two consecutive primes in infinitely many
ways. The claim for non-consecutive primes (which is Polignac's Conjecture as
stated above) is already open.

\subsection{Hardy--Littlewood Generalization}

Hardy and Littlewood (1923) conjectured a quantitative form: for each admissible
even $2k$,
\[
  \#\{p \leq x : p + 2k \text{ prime}\} \sim \mathfrak{S}(2k) \frac{x}{(\log x)^2},
\]
where $\mathfrak{S}(2k) = 2C_2 \prod_{p | k,\, p > 2} \frac{p-1}{p-2}$ is the
\emph{singular series} for the pair $(0, 2k)$.

% ---------------------------------------------------------------------------
\section{Probabilistic Models}
% ---------------------------------------------------------------------------

\subsection{The Cram\'{e}r Model in Detail}

In the Cram\'{e}r model, define a random set $\mathcal{P}$ by: for each $n \geq
3$, include $n$ in $\mathcal{P}$ independently with probability $1/\log n$.
(Adjust $\mathcal{P} \cap [2,3] = \{2,3\}$.)

Under this model:
\begin{enumerate}[label=(\roman*)]
  \item $\mathbb{E}[\pi_{\mathcal{P}}(x)] \sim x/\log x$ (mimics PNT),
  \item gaps near $x$ are approximately exponentially distributed with mean
        $\log x$,
  \item the maximum gap up to $x$ is $\approx (\log x)^2$,
  \item the number of twin prime pairs up to $x$ is
        $\approx x/(\log x)^2$ (but without the $C_2$ correction).
\end{enumerate}

\subsection{The Hardy--Littlewood Heuristic}

A more refined heuristic accounts for the actual density of primes in arithmetic
progressions. For a prime pair $(p, p+2k)$, the local density at $p$ is
$1/\log p$, and the events ``$p$ prime'' and ``$p+2k$ prime'' are not
independent due to common small prime factors. The correction factor for gap $2k$
is the singular series $\mathfrak{S}(2k)$.

For $2k = 2$ (twin primes):
\[
  \mathfrak{S}(2) = 2 C_2 = 2 \prod_{p \geq 3} \frac{p(p-2)}{(p-1)^2} \approx 1.3203.
\]

\subsection{The Granville--Cram\'{e}r Heuristic}

Granville's refined model (1995) predicts that maximal gaps satisfy:
\[
  g_{\max}(x) \approx 2e^{-\gamma} (\log x)^2 \approx 1.1229 (\log x)^2,
\]
where the factor $2e^{-\gamma}$ arises from the asymptotic behavior of the sum
$\sum_{p \leq x} 1/p \approx \log\log x + M$ (Mertens constant $M =
\gamma + \sum_p [\log(1-1/p) + 1/p]$).

% ---------------------------------------------------------------------------
\section{Computational Results}
% ---------------------------------------------------------------------------

\subsection{Twin Prime Data}

The following table shows twin prime counts $\pi_2(x)$ versus the Hardy--Littlewood
prediction $2C_2 \cdot x/(\log x)^2$:

\begin{table}[h]
\centering
\caption{Twin prime counts $\pi_2(x)$ and Hardy--Littlewood prediction}
\begin{tabular}{@{}rrrr@{}}
\toprule
$x$ & $\pi_2(x)$ & $2C_2 x/(\log x)^2$ & Ratio \\
\midrule
$10^3$ & 35 & 27.9 & 1.25 \\
$10^4$ & 205 & 171.3 & 1.20 \\
$10^5$ & 1224 & 1107.8 & 1.10 \\
$10^6$ & 8169 & 7459.7 & 1.10 \\
$10^7$ & 58980 & 54438.1 & 1.08 \\
$10^8$ & 440312 & 412849.5 & 1.07 \\
\bottomrule
\end{tabular}
\end{table}

The ratio approaches $1$ as $x$ increases, consistent with (but not proving)
Conjecture~\ref{conj:hl}.

\subsection{Maximal Gap Records}

\begin{table}[h]
\centering
\caption{Maximal prime gap records (selected)}
\begin{tabular}{@{}rrrr@{}}
\toprule
$p_n$ & $p_{n+1}$ & $g_n$ & $g_n/(\log p_n)^2$ \\
\midrule
7 & 11 & 4 & 0.656 \\
23 & 29 & 6 & 0.594 \\
89 & 97 & 8 & 0.428 \\
113 & 127 & 14 & 0.657 \\
523 & 541 & 18 & 0.503 \\
887 & 907 & 20 & 0.474 \\
1129 & 1151 & 22 & 0.485 \\
9551 & 9587 & 36 & 0.424 \\
31397 & 31469 & 72 & 0.552 \\
155921 & 156007 & 86 & 0.424 \\
360653 & 360749 & 96 & 0.389 \\
\bottomrule
\end{tabular}
\end{table}

\subsection{Andrica Verification}

All Andrica values $A_n = \sqrt{p_{n+1}} - \sqrt{p_n}$ for $n = 1, \ldots,
10^6$ (covering primes up to about $15.5 \times 10^6$) satisfy $A_n < 1$.
The running maximum after $n = 4$ decreases monotonically in recorded maxima:
the global maximum $A_4 \approx 0.6708$ has not been surpassed in any
computation to date.

\subsection{Cram\'{e}r Value Distribution}

For the first $10{,}000$ primes, the normalized gaps $C_n = g_n/(\log p_n)^2$
have:
\begin{itemize}
  \item Mean: $\approx 0.064$,
  \item Maximum: $\approx 0.74$ (near $p = 113$),
  \item The fraction exceeding the Granville constant $2e^{-\gamma} \approx
        1.1229$ is exactly $0$ in this range.
\end{itemize}

These values are consistent with the Cram\'{e}r/Granville heuristics but provide
no proof.

% ---------------------------------------------------------------------------
\section{Summary of Open Problems and Known Results}
% ---------------------------------------------------------------------------

\begin{table}[h]
\centering
\caption{Overview: Open Conjectures and Proven Theorems}
\begin{tabular}{@{}llll@{}}
\toprule
Statement & Author & Year & Status \\
\midrule
Twin Prime Conjecture & Folklore & antiquity & \textbf{Open} \\
Hardy--Littlewood Conj.\ B & Hardy, Littlewood & 1923 & \textbf{Open} \\
Andrica Conjecture & Andrica & 1985 & \textbf{Open} \\
Cram\'{e}r Conjecture & Cram\'{e}r & 1936 & \textbf{Open} \\
Legendre Conjecture & Legendre & $\sim$1798 & \textbf{Open} \\
Polignac Conjecture & de Polignac & 1849 & \textbf{Open} \\
\midrule
Brun's Theorem ($B_2$ converges) & Brun & 1919 & \textbf{Proved} \\
Prime Number Theorem & Hadamard, de la VP & 1896 & \textbf{Proved} \\
$g_n \ll p_n^{0.525}$ & Baker--Harman--Pintz & 2001 & \textbf{Proved} \\
$\liminf g_n \leq 70{,}000{,}000$ & Zhang & 2013 & \textbf{Proved} \\
$\liminf g_n \leq 246$ & Maynard/Polymath8b & 2014 & \textbf{Proved} \\
GPY small gaps & Goldston--Pintz--Y. & 2005 & \textbf{Proved} \\
\bottomrule
\end{tabular}
\end{table}

The central open problems---the Twin Prime Conjecture, Andrica, Cram\'{e}r,
Legendre, and Polignac---all concern the fine structure of prime gaps at various
scales. Zhang's breakthrough shows that the gap $\liminf g_n$ is finite; closing
the enormous distance from $246$ to $2$ remains one of the deepest challenges in
mathematics.

% ---------------------------------------------------------------------------
\section*{References}
% ---------------------------------------------------------------------------

\begin{thebibliography}{99}

\bibitem{brun1919}
V.\ Brun,
\textit{La s\'{e}rie $1/5 + 1/7 + 1/11 + 1/13 + \cdots$ est convergente ou finie},
Bull.\ Sci.\ Math.\ \textbf{43} (1919), 100--104, 124--128.

\bibitem{cramer1936}
H.\ Cram\'{e}r,
\textit{On the order of magnitude of the difference between consecutive prime
numbers},
Acta Arith.\ \textbf{2} (1936), 23--46.

\bibitem{hardylittlewood1923}
G.\ H.\ Hardy and J.\ E.\ Littlewood,
\textit{Some problems of `Partitio Numerorum' III: On the expression of a number
as a sum of primes},
Acta Math.\ \textbf{44} (1923), 1--70.

\bibitem{andrica1985}
D.\ Andrica,
\textit{Note on a conjecture in prime number theory},
Studia Univ.\ Babe\c{s}-Bolyai Math.\ \textbf{31} (1985), no.\ 4, 44--48.

\bibitem{gpy2005}
D.\ A.\ Goldston, J.\ Pintz, and C.\ Y.\ Y{\i}ld{\i}r{\i}m,
\textit{Primes in tuples I},
Ann.\ Math.\ \textbf{170} (2009), 819--862. (Preprint 2005.)

\bibitem{zhang2013}
Y.\ Zhang,
\textit{Bounded gaps between primes},
Ann.\ Math.\ \textbf{179} (2014), 1121--1174. (Submitted 2013.)

\bibitem{maynard2014}
J.\ Maynard,
\textit{Small gaps between primes},
Ann.\ Math.\ \textbf{181} (2015), 383--413. (Preprint 2013.)

\bibitem{polymath8b2014}
D.\ H.\ J.\ Polymath,
\textit{Variants of the Selberg sieve, and bounded intervals containing many
primes},
Res.\ Math.\ Sci.\ \textbf{1} (2014), Art.\ 12.

\bibitem{granville1995}
A.\ Granville,
\textit{Harald Cram\'{e}r and the distribution of prime numbers},
Scand.\ Actuar.\ J.\ \textbf{1} (1995), 12--28.

\bibitem{bhp2001}
R.\ C.\ Baker, G.\ Harman, and J.\ Pintz,
\textit{The difference between consecutive primes, II},
Proc.\ London Math.\ Soc.\ \textbf{83} (2001), 532--562.

\bibitem{westzynthius1931}
E.\ Westzynthius,
\textit{\"{U}ber die Verteilung der Zahlen, die zu den $n$ ersten Primzahlen
teilerfremd sind},
Commentat.\ Phys.-Math.\ \textbf{5} (1931), 1--37.

\end{thebibliography}

\end{document}
